\renewcommand{\thesection}{\arabic{section}}
\chapter{}
\label{chap:p_contenido}


\section{Motivación}

Las aplicaciones web han tenido un crecimiento constante en la última década,
convirtiéndose en infraestructuras críticas para servicios de comercio, comunicación
y administración. Sin embargo, al estar expuestas a través de la red, siguen siendo
objetivo de una amplia variedad de ataques dirigidos al protocolo \gls{acr3:http},
tales como inyecciones, \textit{cross-site scripting} y ataques de traversal
\cite{gimenez2015tfg}.
En muchos casos, las aplicaciones no implementan las mejores prácticas de seguridad,
o bien priorizan la incorporación de funcionalidades por encima de la robustez
defensiva. Esta situación abre la necesidad de nuevas soluciones de detección y
clasificación de ataques que permitan mitigar riesgos de manera efectiva.

En este trabajo nos proponemos expandir los aportes previos sobre detección
de anomalías en tráfico \gls{acr3:http}, integrando mecanismos de
\gls{acr3:ocsvm} con enfoques de clasificación \gls{acr3:ml} multi-etiqueta.
El objetivo es comparar la capacidad de generalización de un clasificador
de una sola clase frente a modelos capaces de identificar explícitamente
el tipo de ataque, aprovechando tanto datasets sintéticos como reales.

Los \gls{acr3:ids} y \gls{acr3:waf} basados en anomalías han demostrado ser
una alternativa valiosa a los enfoques basados en firmas, al no depender de listas
estáticas de ataques conocidos \cite{kruegel2003anomaly}.
No obstante, su aplicación práctica enfrenta el desafío de diferenciar de manera
fiable el tráfico normal de las múltiples variantes de ataques que emergen en la web.
Para abordar esta limitación, nuestro estudio incorpora un nuevo conjunto de
características (\textit{features}) específicas del protocolo \gls{acr3:http},
incluyendo: longitud de la URI, profundidad de la ruta, cantidad de parámetros,
proporción de caracteres no alfanuméricos, porcentaje de codificación hexadecimal,
tokens sospechosos en el contenido y el uso de métodos poco comunes.
La motivación de esta selección radica en que tales rasgos han sido vinculados
directamente con patrones de explotación típicos en ataques web modernos
\cite{nguyen2011application}.

Además de las características tradicionales empleadas en estudios anteriores,
este trabajo incorpora un análisis comparativo entre dos fuentes de datos:
el dataset CSIC 2010/2012 integrado en trabajos como SR-BH 2020,
y un conjunto de tráfico real publicado por Harvard Dataverse \cite{harvard2020dataset}.
Mientras que los primeros son sintéticos y altamente controlados, el segundo
presenta la complejidad y diversidad de un entorno real,
lo que permite contrastar la eficacia de los enfoques en escenarios
con diferentes grados de realismo.

En síntesis, este trabajo busca responder a una necesidad actual:
diseñar mecanismos de detección más robustos y explicativos para el tráfico
\gls{acr3:http}, comparando las ventajas de la detección por una sola clase
con la capacidad clasificadora de los modelos multi-etiqueta.
De esta forma, se espera contribuir tanto a la investigación académica
como al desarrollo práctico de \gls{acr3:waf} inteligentes capaces de
identificar y clasificar ataques de manera automática y precisa.

\section{Justificación}

La problemática de la seguridad en aplicaciones web exige soluciones
que no solo detecten la presencia de tráfico anómalo en el protocolo
\gls{acr3:http}, sino que además permitan caracterizar con precisión
el tipo de ataque.  
Los enfoques basados en firmas resultan limitados, pues dependen
de actualizaciones constantes y no logran cubrir ataques novedosos
o variantes de técnicas existentes \cite{kruegel2003anomaly}.  
Por otro lado, los métodos de detección de anomalías, como
los basados en \gls{acr3:ocsvm}, ofrecen la ventaja de
generalizar frente a comportamientos previamente no observados,
aunque con la desventaja de producir resultados binarios
(normal vs. anómalo).

El aporte de este trabajo consiste en integrar mecanismos de
clasificación multi-etiqueta provenientes del área de
\gls{acr3:ml}, de modo a complementar el poder de generalización
de \gls{acr3:ocsvm} con la capacidad de asignar etiquetas específicas
a los distintos tipos de ataques.  
Esta integración se apoya en la incorporación de un conjunto de
nuevas \textit{features} específicas del protocolo
\gls{acr3:http}, tales como la longitud de la URI, la
profundidad de la ruta, la cantidad de parámetros en la query
y la proporción de caracteres no alfanuméricos.  
Estos rasgos han demostrado ser indicadores relevantes en la
detección de anomalías y en la diferenciación de intentos de
explotación \cite{nguyen2011application}.

Otro aspecto que justifica este estudio es la comparación entre
datasets de distinta naturaleza.  
Mientras que conjuntos como CSIC 2010/2012 se construyeron en
entornos controlados, el dataset de Harvard Dataverse refleja
el comportamiento real del tráfico en un ambiente productivo
\cite{harvard2020dataset}.  
Esta dualidad permitirá evaluar de manera más completa la
eficacia y robustez de los clasificadores bajo escenarios
diferentes.

En consecuencia, el trabajo se justifica en tres dimensiones:
(1) relevancia práctica en la protección de aplicaciones web
modernas, (2) aporte académico al comparar \gls{acr3:ocsvm}
y clasificación multi-etiqueta en tráfico \gls{acr3:http},
y (3) innovación en el diseño y evaluación de nuevas
\textit{features} que enriquecen la detección de ataques.

\section{Objetivos}

\subsection{Objetivo general}

Comparar la eficacia de un clasificador \gls{acr3:ocsvm} y de
enfoques de clasificación multi-etiqueta para la detección y
caracterización de ataques en tráfico \gls{acr3:http}, integrando
nuevas \textit{features} y evaluando el desempeño en datasets
sintéticos (CSIC 2010/2012) y reales (Harvard Dataverse).

\subsection{Objetivos específicos}

\begin{enumerate}
  \item Diseñar y extraer un conjunto ampliado de
        \textit{features} específicas del protocolo \gls{acr3:http},
        incorporando métricas como longitud de la URI, profundidad
        de la ruta, cantidad de parámetros y tokens sospechosos.

  \item Implementar y entrenar un clasificador
        \gls{acr3:ocsvm} para la detección de tráfico anómalo,
        utilizando los datasets seleccionados.

  \item Implementar clasificadores multi-etiqueta basados en
        \gls{acr3:ml}, capaces de asignar a cada muestra
        el tipo específico de ataque.

  \item Comparar los resultados de ambos enfoques en términos
        de precisión, \textit{recall}, F1-score y tasa de falsos
        positivos/negativos, considerando tanto el dataset CSIC
        como el de Harvard Dataverse.

  \item Analizar el impacto de las nuevas \textit{features}
        en la capacidad de detección y clasificación de los modelos,
        identificando cuáles contribuyen más a la discriminación
        entre tráfico normal y malicioso.

  \item Evaluar la viabilidad de aplicar los enfoques
        propuestos en escenarios de tráfico real en tiempo
        casi real, considerando el costo computacional y la
        escalabilidad.
\end{enumerate}

\section{Descripción de la propuesta - aca todavia no esta}

Nuestra propuesta en el marco de este trabajo consiste en un \gls{acr3:waf}
que puede ser colocado frente a varias aplicaciones web con el fin de
monitorear todo el tráfico \gls{acr3:http} entre dichas aplicaciones y
sus usuarios. En la \autoref{fig:arquitectura} se puede observar la
arquitectura general de nuestra propuesta.

\begin{figure}[ht]
    \centering
    \includegraphics[width=\linewidth]{images/waf-diagram.png}
    \caption{Diagrama de la arquitectura general de la propuesta}
    \label{fig:arquitectura}
\end{figure}

La implementación tendrá la opción de habilitar o deshabilitar la detección
en tiempo real. Con la detección deshabilitada, el \gls{acr3:waf} se
limitará a registrar el tráfico para que posteriormente se pueda realizar
el entrenamiento del clasificador con muestras normales de tráfico.
Una vez entrenado, se podrá habilitar el modo de detección. En este modo
el \gls{acr3:waf} analizará todas las peticiones entrantes y registrará
todas aquellas que sean clasificadas como anómalas.
Se podrá configurar acciones adicionales que el \gls{acr3:waf} deberá
realizar como respuesta a una anomalía detectada, siendo una opción
el bloqueo de la petición en cuestión para evitar que posibles ataques
lleguen hasta las aplicaciones.


\section{Alcance y limitaciones}

Según nuestro relevamiento inicial, el trabajo de Kruegel
\cite{kruegel2003anomaly} es la investigación pionera relacionada con
\gls{acr3:ids} que se enfoca en las particularidades del tráfico HTTP,
especialmente el ánalisis de los parámetros y valores dentro de la
petición. Con más de 600 citas hasta la fecha según Google Scholar,
este trabajo ha sido utilizado como base para muchas investigaciones
posteriores.

De esta manera, la investigación del estado del arte que haremos en el
marco de nuestro trabajo se limitará a las publicaciones que citan a Kruegel
en sus referencias.
\bigskip

Para la evaluación cuantitativa de nuestro \gls{acr3:waf} propuesto
necesitamos conjuntos de datos. Según \cite{torranoGimenez2015study}
los conjuntos utilizados en las investigaciones sobre \gls{acr3:waf}
deberían cumplir las siguientes características:

\begin{itemize}
    \item
    Deberían ser públicamente accesibles para que varios investigadores
    los utilizen y puedan comparar resultados.

    \item
    Deberían contener tráfico \gls{acr3:http}.

    \item
    Deberían contener muestras etiquetadas según sean normales o
    anómalas, o inclusive especificar el tipo de ataque al que
    pertenecen.

    \item
    Deberían tener ataques novedosos.
\end{itemize}

Los trabajos relacionados sobre \gls{acr3:waf} utilizan distintos
conjuntos de datos para sus pruebas, pero la mayoria de esos
conjuntos no cumple una o varias de las características mencionadas.
Encontramos unicamente dos conjuntos que son adecuados para
nuestras pruebas y se trata de CSIC 2010 \cite{csic2010dataset} y
CSIC TORPEDA 2012 \cite{torpeda2012dataset}. Así podemos realizar
comparaciones con los resultados de otros trabajos que utilizan estos
conjuntos.
\bigskip

La implementación de un \gls{acr3:waf} que haremos en el marco de este
trabajo busca ser funcional y sencilla. Se trata de una prueba de concepto,
y no buscamos obtener una aplicación terminada que incluya todas las
funcionalidades necesarias para utilizarla directamente en ambientes de
producción.


\section{Plan de actividades}

A continuación se detallan las actividades que nosotros nos proponemos
para el desarrollo de este trabajo de investigación. Además, en la
\autoref{tbl:actividades} se puede observar las duraciones estimadas
para dichas actividades.

\begin{enumerate}
    \item
    Exploración bibliográfica sobre las áreas de \gls{acr3:ids},
    \gls{acr3:waf} y \gls{acr3:occ}, profundizando especialmente
    en trabajos que utilizan el clasificador \gls{acr3:ocsvm} y
    aquellos que basan sus \textit{features} en los trabajos de Kruegel.

    \item
    Diseño de nuevos conjuntos de \textit{features} específicos para
    tráfico \gls{acr3:http}, combinando las ideas de Kruegel con aportes
    propios y también de otros investigadores.

    \item
    Construcción de un \gls{acr3:waf} con los \textit{features} diseñados
    y con un clasificador \gls{acr3:ocsvm}.

    \item
    Pruebas de eficacia de detección del \gls{acr3:waf} construido,
    buscando conjuntos de \textit{features} y configuraciones del
    clasificador que mejoren la detección.

    \item
    Pruebas de rendimiento del \gls{acr3:waf} construido, analizando
    en qué medida el mismo afecta al tiempo de respuesta de los
    sistemas que debe proteger.

    \item
    Análisis de resultados y formulación de conclusiones.
\end{enumerate}

\begin{table}[ht]
    \centering
    \small
    \begin{tabular*}{\linewidth}{@{\extracolsep{\fill}}|l|c|c|c|c|c|c|c|c|c|c|c|c|c|c|c|c|c|c|}
        \hline
        Año & \multicolumn{18}{c|}{2017} \\ \hline
        Mes & \multicolumn{5}{c|}{Agosto} & \multicolumn{4}{c|}{Setiembre} & \multicolumn{4}{c|}{Octubre} & \multicolumn{5}{c|}{Noviembre} \\ \hline
        Semana & 1 & 2 & 3 & 4 & 5 & 1 & 2 & 3 & 4 & 1 & 2 & 3 & 4 & 1 & 2 & 3 & 4 & 5 \\ \hline \hline
        Actividad 1 & x & x & x & x & x & x & x &  &  &  &  &  &  &  &  &  &  &  \\ \hline
        Actividad 2 &  &  &  &  & x & x & x &  &  &  &  &  &  &  &  &  &  &  \\ \hline
        Actividad 3 &  &  &  &  &  &  &  & x & x & x & x &  &  &  &  &  &  &  \\ \hline
        Actividad 4 &  &  &  &  &  &  &  &  &  &  &  & x & x & x &  &  &  &  \\ \hline
        Actividad 5 &  &  &  &  &  &  &  &  &  &  &  & x & x & x &  &  &  &  \\ \hline
        Actividad 6 &  &  &  &  &  &  &  &  &  &  &  &  &  &  & x & x & x & x \\ \hline
    \end{tabular*}
    \caption{Duración estimada de actividades}
    \label{tbl:actividades}
\end{table}
